%---------------------- Problema ---------------------------------
\section*{Problema}

\begin{enunciado}{Problema (6 puntos)}
    \begin{enumerate}[label=\alph*), itemsep=4pt]
        \item Calcular mediante el m\'etodo de elementos finitos los campos de
        las diferentes variables correspondientes al problema del apartado 35
        del libro \textit{``Theory of Elasticity''} de S.P.~Timoshenko y
        J.N.~Goodier (todos los par\'ametros son de libre elecci\'on):
        \textit{``The Effect of Circular Holes on Stress Distributions in
        Plates''} (figura 52).
        \item Comparar y discutir los resultados obtenidos en a) con los
        correspondientes a la soluci\'on anal\'itica del mencionado apartado.
        \item Comparar y discutir los resultados obtenidos en a) con el
        proporcionado por un programa comercial de elementos finitos.
    \end{enumerate}
\end{enunciado}

\begin{figure}[!htbp]
\centering
\begin{tikzpicture}[>=latex, scale=1.1]
    % --- Datos geometricos (solo para el dibujo) ---
    \def\Lplaca{9}      % largo de la placa
    \def\dplaca{3}      % ancho de la placa (d)
    \def\rag{0.6}    % radio del agujero (r)
    \pgfmathsetmacro{\xc}{\Lplaca/2}
    \pgfmathsetmacro{\yc}{\dplaca/2}

    % placa con agujero
    \draw[thick, fill=black!5, even odd rule]
        (0,0) rectangle (\Lplaca,\dplaca)
        (\xc,\yc) circle (\rag);

    % ejes de simetria
    \draw[dashed, thin, gray] (-0.2,\yc) -- (\Lplaca+0.2,\yc);
    \draw[dashed, thin, gray] (\xc,-0.3) -- (\xc,\dplaca+0.3);

    % traccion uniforme S en ambos extremos
    \foreach \y in {0.25,0.75,...,2.75}{
        \draw[-{Latex[length=2mm]}, thick] (0,\y) -- (-0.8,\y);
        \draw[-{Latex[length=2mm]}, thick] (\Lplaca,\y) -- (\Lplaca+0.8,\y);
    }
    \node at (-1.2,\yc) {\Large $S$};
    \node at (\Lplaca+1.2,\yc) {\Large $S$};

    % puntos m (borde de la placa) y n (borde del agujero)
    \fill (\xc,\dplaca) circle (1.6pt) node[above=3pt] {$m$};
    \fill (\xc,\yc+\rag) circle (1.6pt) node[above left=1pt] {$n$};

    % cota 2r
    \draw[<->, thin] (\xc+1.1,\yc-\rag) -- (\xc+1.1,\yc+\rag)
        node[midway, fill=black!5, inner sep=2pt] {$2r$};
    \draw[thin] (\xc,\yc-\rag) -- (\xc+1.25,\yc-\rag);
    \draw[thin] (\xc,\yc+\rag) -- (\xc+1.25,\yc+\rag);

    % cota d
    \draw[<->, thick] (\Lplaca+2.0,0) -- (\Lplaca+2.0,\dplaca)
        node[midway, fill=white, inner sep=2pt] {$d$};
    \draw[thin] (\Lplaca+0.1,0) -- (\Lplaca+2.2,0);
    \draw[thin] (\Lplaca+0.1,\dplaca) -- (\Lplaca+2.2,\dplaca);
\end{tikzpicture}
\caption{Placa de ancho $d$ con un agujero circular de di\'ametro $2r$ en su
eje de simetr\'ia, sometida a tracci\'on uniforme $S$ en sus extremos
(Timoshenko \& Goodier, Fig.~52).}
\label{fig:t3-placa}
\end{figure}

%---------------------- Desarrollo -------------------------------

\begin{desarrollo}{a) Soluci\'on por MEF de los campos de la placa con agujero}
    % Par\'ametros elegidos (S, r, d, largo, espesor, E, nu), modelo (tensi\'on
    % plana / deformaci\'on plana, simetr\'ia), elemento, malla, condiciones de
    % borde, c\'odigo y resultados (desplazamientos, sigma_x, sigma_y, tau_xy,
    % sigma_r, sigma_theta, tau_rtheta).

\end{desarrollo}

\begin{desarrollo}{b) Comparaci\'on con la soluci\'on anal\'itica}
    % Soluci\'on de Kirsch (Ec. 61), sigma_theta en el borde del agujero,
    % sigma_theta en theta = pi/2 vs r, factor de concentraci\'on en n y m,
    % efecto del ancho finito d.

\end{desarrollo}

\begin{desarrollo}{c) Comparaci\'on con un programa comercial de elementos finitos}
    % Programa usado (ANSYS, Abaqus, SolidWorks, ...), modelo, malla,
    % resultados y comparaci\'on con a) y b).

\end{desarrollo}
